% Capítulo — TCC 1
\chapter{Introdução}
\label{chap:introducao}

% Abertura do capítulo: um parágrafo dizendo o que ele apresenta.

\section{Contexto}
\label{sec:introducao-contexto}
% Dados públicos federais em silos; o GovHub-br e a parceria UnB--ministérios.

\section{Problema}
\label{sec:introducao-problema}
% O de-para entre bases (códigos, nomes e formatos diferentes) é manual e depende de conhecimento implícito do analista.

\section{Justificativa}
\label{sec:introducao-justificativa}
% Custo de integrar cada nova fonte; o padrão do SPAPI-Tester (LLM embutido no processo) ainda não aplicado a dados abertos.

\section{Questões de pesquisa}
\label{sec:introducao-questoes-de-pesquisa}
% QP1 a QP5; tabela com pergunta e como será respondida.

\section{Objetivos}
\label{sec:introducao-objetivos}

\subsection{Objetivo geral}
\label{subsec:introducao-objetivo-geral}
% Propor, implementar e avaliar uma arquitetura config-driven de três camadas que integre bases governamentais heterogêneas, com LLM embutido para descobrir chaves de integração.

\subsection{Objetivos específicos}
\label{subsec:introducao-objetivos-especificos}
% Lista numerada: Source Registry, costuras, Decision Layer com DSPy, ground truth, avaliação, construção por SDD (QP5).

\section{Contribuições esperadas}
\label{sec:introducao-contribuicoes-esperadas}
% Arquitetura aberta, gerador dbt config-driven, benchmark de chaves por categoria de atrito, replicação do artigo-base, relato avaliado de SDD.

\section{Organização do trabalho}
\label{sec:introducao-organizacao-do-trabalho}
% Um parágrafo por capítulo.
